\documentclass{article}
\usepackage{graphicx} % Required for inserting images


\usepackage[toc, sort=def, nonumberlist]{glossaries}
\makeglossaries


\newglossaryentry{mdp}{
    name={MDP},
    description={Markov Decision Process. A mathematical framework used to model decision-making situations where outcomes are partly random and partly under the control of a decision-maker}
}

\newglossaryentry{state}{
    name={State Space ($S$)},
    description={The set of all valid situations, configurations, or positions ($s \in S$) that the agent can occupy within the environment}
}

\newglossaryentry{action}{
    name={Action Space ($A$)},
    description={The set of all possible moves, choices, or decisions ($a \in A$) available to the agent in its current situation}
}

\newglossaryentry{transition}{
    name={Transition Probability ($P$)},
    description={The probability $P(s' \mid s, a)$ that executing a specific action $a$ in state $s$ will cause the environment to shift to a specific next state $s'$. This accounts for natural randomness or noise in the environment}
}

\newglossaryentry{reward}{
    name={Reward Function ($R$)},
    description={The immediate numerical feedback $R(s, a, s')$ given to the agent by the environment right after moving from state $s$ to state $s'$ via action $a$}
}

\newglossaryentry{policy}{
    name={Policy ($\pi$)},
    description={The agent's strategic brain or rulebook. Formally written as $\pi(a \mid s)$, it maps a given state to the probability distribution of selecting each available action. The ultimate goal of the agent is to find the Optimal Policy ($\pi^*$)}
}

\newglossaryentry{statevalue}{
    name={State-Value Function ($V$)},
    description={Denoted as $V(s)$, it measures how good it is to be in a specific state $s$ by calculating the total expected future reward accumulated from that point onward, assuming the agent follows its current strategy until the end}
}

\newglossaryentry{actionvalue}{
    name={Action-Value Function ($Q$)},
    description={The Quality function, denoted as $Q(s, a)$, which calculates the total expected long-term return if the agent commits to a specific initial action $a$ from state $s$, and follows its policy thereafter}
}

\newglossaryentry{discount}{
    name=Discount Factor,
    description={A factor $\gamma \in [0,1]$ that determines the present value of future rewards}
}
\newglossaryentry{transitionprobability}{
    name=Transition Probability,
    description={The probability $P(s'|s,a)$ of transitioning to state $s'$ from state $s$ after taking action $a$}
}
\newglossaryentry{rewardfunction}{
    name=Reward Function,
    description={The function $R(s,a,s')$ that gives the immediate reward for transitioning from state $s$ to $s'$ via action $a$}
}


\title{Week 2}
\author{Prerit Advani}
\date{June 2026}

\begin{document}

\maketitle

\section{Solving MDPs through Model-Free Methods (Reinforcement Learning)}

When the environment is perfectly modeled, we possess exact knowledge of the \gls{transitionprobability} $P(s'|s, a)$ and the \gls{rewardfunction} $R(s, a, s')$. In this theoretical setting, we can use \textbf{Dynamic Programming (DP)} to solve for the \gls{statevalue} or \gls{actionvalue} functions. By iteratively applying the Bellman equations, we can compute the \gls{policy} that yields the highest possible return without ever needing to interact with the environment.


But, In most real-world scenarios, the transition probabilities and reward functions are unknown. Here, we use \textbf{Reinforcement Learning (RL)} to estimate these values through trial and error.

\subsection{temporal Differencing (TD)}

Methods that learn by observing the differences between predicted rewards and actual rewards received after taking an action. There are primarily 2 temporal differencing methods used:

\subsubsection{Q-Learning (Off-Policy)}

The agent learns the value of the optimal policy regardless of the actions it is actually taking. It updates its $Q(s, a)$ estimates based on the maximum possible future reward.

\subsubsection{SARSA (On-policy)}

The agent learns the value of the policy it is currently following, including the exploratory moves it makes. It updates $Q(s, a)$ based on the action it actually takes next.

\subsection{Monte Carlo Methods}

These learn from complete episodes of experience by averaging the total returns received after visiting a state. Unlike TD learning, they only update after an episode ends.

\subsection{Deep Actor-Critic Methods}
While traditional TD methods (like Q-Learning) learn value functions to derive a policy, modern Deep RL often employs \textbf{Actor-Critic} architectures. These methods combine the benefits of policy-based learning (learning the strategy directly) and value-based learning (using a \textit{critic} to evaluate actions via TD errors). Examples include:

\begin{itemize}
    \item \textbf{PPO (Proximal Policy Optimization):} An on-policy method focused on stable, incremental policy updates.
    \item \textbf{DDPG (Deep Deterministic Policy Gradient):} An off-policy method designed for continuous action spaces.
    \item \textbf{SAC (Soft Actor-Critic):} An off-policy method that maximizes both reward and entropy for robust exploration.
\end{itemize}

\newpage
\glsaddall
\printglossaries

\end{document}
